\section{More on invertible modules}
\label{section-invertible-modules}

\noindent
In this section we discuss some properties of invertible modules.

\begin{lemma}
\label{lemma-in-image-pullback}
Let $\varphi : X \to Y$ be a morphism of schemes. Let $\mathcal{L}$
be an invertible $\mathcal{O}_X$-module. Assume that
\begin{enumerate}
\item $X$ is locally Noetherian,
\item $Y$ is locally Noetherian, integral, and normal,
\item $\varphi$ is flat with integral (hence nonempty) fibres,
\item $\varphi$ is either quasi-compact or locally of finite type,
\item $\mathcal{L}$ is trivial when restricted to the generic fibre of
$\varphi$.
\end{enumerate}
Then $\mathcal{L} \cong \varphi^*\mathcal{N}$ for some invertible
$\mathcal{O}_Y$-module $\mathcal{N}$.
\end{lemma}

\begin{proof}
Let $\xi \in Y$ be the generic point. Let $X_\xi$ be the scheme theoretic
fibre of $\varphi$ over $\xi$. Denote $\mathcal{L}_\xi$ the pullback of
$\mathcal{L}$ to $X_\xi$. Assumption (5) means that $\mathcal{L}_\xi$
is trivial. Choose a trivializing section
$s \in \Gamma(X_\xi, \mathcal{L}_\xi)$. Observe that $X$ is integral by
Lemma \ref{lemma-flat-over-integral-integral-fibre}.
Hence we can think of $s$ as a regular meromorphic section of $\mathcal{L}$.
Pullbacks of meromorphic functions are defined for
$\varphi$ by Lemma \ref{lemma-pullback-meromorphic-sections-defined}.
Let $\mathcal{N} \subset \mathcal{K}_Y$ be the $\mathcal{O}_Y$-module
whose sections over an open $V \subset Y$ are those meromorphic functions
$g \in \mathcal{K}_Y(V)$ such that
$\varphi^*(g)s \in \mathcal{L}(\varphi^{-1}V)$.
A priori $\varphi^*(g)s$ is a section of $\mathcal{K}_X(\mathcal{L})$ over
$\varphi^{-1}V$. We claim that $\mathcal{N}$ is an invertible
$\mathcal{O}_Y$-module and that the map
$$
\varphi^*\mathcal{N} \longrightarrow \mathcal{L},\quad
g \longmapsto gs
$$
is an isomorphism.

\medskip\noindent
We first prove the claim in the following situation:
$X$ and $Y$ are affine and $\mathcal{L}$ trivial. Say $Y = \Spec(R)$,
$X = \Spec(A)$ and $s$ given by the element $s \in A \otimes_R K$
where $K$ is the fraction field of $R$. We can write $s = a/r$
for some nonzero $r \in R$ and $a \in A$. Since $s$ generates $\mathcal{L}$
on the generic fibre we see that there exists an $s' \in A \otimes_R K$
such that $ss' = 1$. Thus we see that $s = r'/a'$ for some
nonzero $r' \in R$ and $a' \in A$. Let
$\mathfrak p_1, \ldots, \mathfrak p_n \subset R$
be the minimal primes over $rr'$. Each $R_{\mathfrak p_i}$
is a discrete valuation ring
(Algebra, Lemmas \ref{algebra-lemma-minimal-over-1} and
\ref{algebra-lemma-criterion-normal}). By assumption
$\mathfrak q_i = \mathfrak p_i A$ is a prime. Hence
$\mathfrak q_i A_{\mathfrak q_i}$ is generated by a single element
and we find that $A_{\mathfrak q_i}$ is a discrete valuation ring as
well (Algebra, Lemma \ref{algebra-lemma-characterize-dvr}). Of course
$R_{\mathfrak p_i} \to A_{\mathfrak q_i}$ has ramification index $1$.
Let $e_i, e'_i \geq 0$ be the  valuation of $a, a'$ in $A_{\mathfrak q_i}$.
Then $e_i + e'_i$ is the valuation of $rr'$ in $R_{\mathfrak p_i}$. Note that
$$
\mathfrak p_1^{(e_1 + e'_1)} \cap \ldots \cap \mathfrak p_i^{(e_n + e'_n)} =
(rr')
$$
in $R$ by Algebra, Lemma
\ref{algebra-lemma-normal-domain-intersection-localizations-height-1}.
Set
$$
I = \mathfrak p_1^{(e_1)} \cap \ldots \cap \mathfrak p_i^{(e_n)}
\quad\text{and}\quad
I' = \mathfrak p_1^{(e'_1)} \cap \ldots \cap \mathfrak p_i^{(e'_n)}
$$
so that $II' \subset (rr')$. Observe that
$$
IA =
(\mathfrak p_1^{(e_1)} \cap \ldots \cap \mathfrak p_i^{(e_n)})A =
(\mathfrak p_1A)^{(e_1)} \cap \ldots \cap (\mathfrak p_i A)^{(e_n)}
$$
by Algebra, Lemmas \ref{algebra-lemma-symbolic-power-flat-extension} and
\ref{algebra-lemma-flat-intersect-ideals}. Similarly for $I'A$. Hence
$a \in IA$ and $a' \in I'A$.
We conclude that $IA \otimes_A I'A \to rr'A$ is surjective.
By faithful flatness of $R \to A$ we find that
$I \otimes_R I' \to (rr')$ is surjective as well.
It follows that $II' = (rr')$ and $I$ and $I'$ are finite locally
free of rank $1$, see
Algebra, Lemma \ref{algebra-lemma-product-ideals-principal}.
Thus Zariski locally on $R$ we can write $I = (g)$ and $I' = (g')$
with $gg' = rr'$. Then $a = ug$ and $a' = u'g'$ for some $u, u' \in A$.
We conclude that $u, u'$ are units. Thus Zariski locally on $R$
we have $s = ug/r$ and the claim follows in this case.

\medskip\noindent
Let $y \in Y$ be a point.
Pick $x \in X$ mapping to $y$. We may apply the result of the previous
paragraph to $\Spec(\mathcal{O}_{X, x}) \to \Spec(\mathcal{O}_{Y, y})$.
We conclude there exists an element $g \in R(Y)^*$ well defined up to
multiplication by an element of $\mathcal{O}_{Y, y}^*$ such that
$\varphi^*(g)s$ generates $\mathcal{L}_x$. Hence $\varphi^*(g)s$
generates $\mathcal{L}$ in a neighbourhood $U$ of $x$.
Suppose $x'$ is a second point lying over $y$ and $g' \in R(Y)^*$ is
such that $\varphi^*(g')s$ generates $\mathcal{L}$ in an open neighbourhood
$U'$ of $x'$. Then we can choose a point
$x''$ in $U \cap U' \cap \varphi^{-1}(\{y\})$
because the fibre is irreducible. By the uniqueness for
the ring map $\mathcal{O}_{Y, y} \to \mathcal{O}_{X, x''}$
we find that $g$ and $g'$ differ (multiplicatively)
by an element in $\mathcal{O}_{Y, y}^*$. Hence we see that $\varphi^*(g)s$
is a generator for $\mathcal{L}$ on an open neighbourhood
of $\varphi^{-1}(y)$. Let $Z \subset X$ be the set of points
$z \in X$ such that $\varphi^*(g)s$ does not generate $\mathcal{L}_z$.
The arguments above show that $Z$ is closed and that $Z = \varphi^{-1}(T)$
for some subset $T \subset Y$ with $y \not \in T$. If we can show that
$T$ is closed, then $g$ will be a generator for $\mathcal{N}$ as an
$\mathcal{O}_Y$-module in the open neighbourhood $Y \setminus T$ of $y$
thereby finishing the proof (some details omitted).

\medskip\noindent
If $\varphi$ is quasi-compact, then $T$ is closed by
Morphisms, Lemma \ref{morphisms-lemma-fpqc-quotient-topology}.
If $\varphi$ is locally of finite type, then $\varphi$ is open
by Morphisms, Lemma \ref{morphisms-lemma-fppf-open}.
Then $Y \setminus T$ is open as the image of the open $X \setminus Z$.
\end{proof}

\begin{lemma}
\label{lemma-closure-effective-cartier-divisor}
Let $X$ be a locally Noetherian scheme. Let $U \subset X$ be
an open and let $D \subset U$ be an effective Cartier divisor.
If $\mathcal{O}_{X, x}$ is a UFD for all $x \in X \setminus U$,
then there exists an effective Cartier divisor $D' \subset X$
with $D = U \cap D'$.
\end{lemma}

\begin{proof}
Let $D' \subset X$ be the scheme theoretic image of the morphism $D \to X$.
Since $X$ is locally Noetherian the morphism $D \to X$ is quasi-compact, see
Properties, Lemma \ref{properties-lemma-immersion-into-noetherian}.
Hence the formation of $D'$ commutes with passing to opens in $X$ by
Morphisms, Lemma \ref{morphisms-lemma-quasi-compact-scheme-theoretic-image}.
Thus we may assume $X = \Spec(A)$ is affine.
Let $I \subset A$ be the ideal corresponding to $D'$.
Let $\mathfrak p \subset A$ be a prime ideal corresponding to a point
of $X \setminus U$.
To finish the proof it is enough to show that $I_\mathfrak p$ is generated
by one element, see Lemma \ref{lemma-effective-Cartier-in-points}.
Thus we may replace $X$ by $\Spec(A_\mathfrak p)$, see
Morphisms, Lemma \ref{morphisms-lemma-flat-base-change-scheme-theoretic-image}.
In other words, we may assume that $X$ is the spectrum of a local
UFD $A$. Then all local rings of $A$ are UFD's. It follows that
$D = \sum a_i D_i$ with $D_i \subset U$ an integral effective Cartier divisor,
see Lemma \ref{lemma-effective-Cartier-divisor-is-a-sum}.
The generic points $\xi_i$ of $D_i$ correspond to prime ideals
$\mathfrak p_i \subset A$ of height $1$, see
Lemma \ref{lemma-effective-Cartier-codimension-1}.
Then $\mathfrak p_i = (f_i)$ for some prime element $f_i \in A$
and we conclude that $D'$ is cut out by $\prod f_i^{a_i}$ as desired.
\end{proof}

\begin{lemma}
\label{lemma-extend-invertible-module}
Let $X$ be a locally Noetherian scheme. Let $U \subset X$ be
an open and let $\mathcal{L}$ be an invertible $\mathcal{O}_U$-module.
If $\mathcal{O}_{X, x}$ is a UFD for all $x \in X \setminus U$,
then there exists an invertible $\mathcal{O}_X$-module $\mathcal{L}'$
with $\mathcal{L} \cong \mathcal{L}'|_U$.
\end{lemma}

\begin{proof}
Choose $x \in X$, $x \not \in U$. We will show there exists an
affine open neighbourhood $W \subset X$, such that $\mathcal{L}|_{W \cap U}$
extends to an invertible sheaf on $W$. This implies by glueing of
sheaves (Sheaves, Section \ref{sheaves-section-glueing-sheaves})
that we can extend $\mathcal{L}$ to the strictly bigger open $U \cup W$.
Let $W = \Spec(A)$ be an affine open neighbourhood.
Since $U \cap W$ is quasi-affine, we see that we can write
$\mathcal{L}|_{W \cap U}$ as
$\mathcal{O}(D_1) \otimes \mathcal{O}(D_2)^{\otimes -1}$ for some
effective Cartier divisors $D_1, D_2 \subset W \cap U$, see
Lemma \ref{lemma-quasi-projective-Noetherian-pic-effective-Cartier}.
Then $D_1$ and $D_2$ extend to effective Cartier divisors of
$W$ by Lemma \ref{lemma-closure-effective-cartier-divisor}
which gives us the extension of the invertible sheaf.

\medskip\noindent
If $X$ is Noetherian (which is the case most used in practice), the above
combined with Noetherian induction finishes the proof. In the general case
we argue as follows. First, because every local ring of a point outside
of $U$ is a domain and $X$ is locally Noetherian, we see that the closure
of $U$ in $X$ is open. Thus we may assume that $U \subset X$ is dense
and schematically dense.
Now we consider the set $T$ of triples $(U', \mathcal{L}', \alpha)$
where $U \subset U' \subset X$ is an open subscheme, $\mathcal{L}'$
is an invertible $\mathcal{O}_{U'}$-module, and
$\alpha : \mathcal{L}'|_U \to \mathcal{L}$ is an isomorphism.
We endow $T$ with a partial ordering $\leq$ defined by the rule
$(U', \mathcal{L}', \alpha) \leq (U'', \mathcal{L}'', \alpha')$
if and only if $U' \subset U''$ and there exists an isomorphism
$\beta : \mathcal{L}''|_{U'} \to \mathcal{L}'$ compatible with
$\alpha$ and $\alpha'$. Observe that $\beta$ is unique (if it exists)
because $U \subset X$ is dense. The first part of the proof shows that
for any element $t = (U', \mathcal{L}', \alpha)$ of $T$ with $U' \not = X$
there exists a $t' \in T$ with $t' > t$. Hence to finish the proof it
suffices to show that Zorn's lemma applies. Thus consider a
totally ordered subset $I \subset T$. If $i \in I$ corresponds to
the triple $(U_i, \mathcal{L}_i, \alpha_i)$, then we can construct
an invertible module $\mathcal{L}'$ on $U' = \bigcup U_i$ as follows.
For $W \subset U'$ open and quasi-compact we see that
$W \subset U_i$ for some $i$ and we set
$$
\mathcal{L}'(W) = \mathcal{L}_i(W)
$$
For the transition maps we use the $\beta$'s (which are unique and hence
compose correctly). This defines an invertible $\mathcal{O}$-module
$\mathcal{L}'$ on the basis of quasi-compact opens of $U'$ which is sufficient
to define an invertible module (Sheaves, Section \ref{sheaves-section-bases}).
We omit the details.
\end{proof}

\begin{lemma}
\label{lemma-open-subscheme-UFD}
Let $R$ be a UFD. The Picard groups of the following are
trivial.
\begin{enumerate}
\item $\Spec(R)$ and any open subscheme of it.
\item $\mathbf{A}^n_R = \Spec(R[x_1, \ldots, x_n])$ and any open subscheme
of it.
\end{enumerate}
In particular, the Picard group of any open subscheme of affine
$n$-space $\mathbf{A}^n_k$ over a field $k$ is trivial.
\end{lemma}

\begin{proof}
Any polynomial ring over $R$ is a UFD by
Algebra, Lemma \ref{algebra-lemma-polynomial-ring-UFD}.
Thus the question is to show that $\Pic(U)$ is trivial
if $U \subset \Spec(R)$ is a (nonempty) open.

\medskip\noindent
In the Noetherian case, we may use Lemma \ref{lemma-extend-invertible-module}
to reduce to showing $\Pic(\Spec(R))$ is trivial. This in turn is
proved in More on Algebra, Lemma \ref{more-algebra-lemma-UFD-Pic-trivial}.
We urge the reader to skip the rest of the proof.

\medskip\noindent
Let $U \subset \Spec(R)$ be a nonempty open. Choose a nonzero $f \in R$
with $D(f) \subset U$. Since $f$ is a product of prime elements of $R$
we conclude that there are finitely many prime elements
$a_1, \ldots, a_n \in R$ such that $(a_i) \not \in U$ (or
equivalentaly $U \cap V(a_i) = \emptyset$).
After replacing $R$ by $R_{a_1 \ldots a_n}$ we may assume that
$U$ contains all principal prime ideals. (Note that a localization
of a UFD is a UFD, for example by Algebra, Lemma
\ref{algebra-lemma-invert-prime-elements}.)

\medskip\noindent
Assume $U$ contains all principal prime ideals.
Write $U = \Spec(R) \setminus V(I)$ for some ideal $I \subset R$.
Pick $f \in I$ nonzero and write $f = a_1 \ldots a_n$ as a product
of prime elements. Since $(a_i) \not \in V(I)$ we may choose $g \in I$ with
$g \not \in (a_i)$ for all $i$ (prime avoidance). Picture
$$
U' = D(f) \cup D(g) \subset U \subset \Spec(R)
$$
By More on Algebra, Lemma \ref{more-algebra-lemma-UFD-Pic-trivial}
the Picard groups of $D(f)$ and $D(g)$ are trivial and by an elementary
argument about UFDs the map $R_f^* \times R_g^* \to R_{fg}^*$,
$(u, v) \mapsto uv^{-1}$ is surjective. We conclude that
$\Pic(U')$ is trivial.

\medskip\noindent
Thus it suffices to show that an invertible module $\mathcal{L}$
on $U$ whose pullback by $j : U' \to U'$ is trivial, is trivial.
To do this it suffices to show that the map
$\mathcal{L} \to j_*j^*\mathcal{L}$ is an isomorphism.
For this in turn, let $u \in U \setminus U'$. Then
$f, g \in \mathfrak m_u \subset \mathcal{O}_{U, u} = A$ are elements
of a local UFD which do not have a prime factor in common. The reader
easily shows that this means that
$$
0 \to A \to A_f \times A_g \to A_{fg}
$$
is exact by looking at prime factors. Since $A = \mathcal{L}_u$
because $\mathcal{L}$ is invertible and since $(j_*j^*\mathcal{L})_u$
is the kernel of $A_f \times A_g \to A_{fg}$ we conclude.
\end{proof}

\begin{lemma}
\label{lemma-Pic-projective-space-UFD}
Let $R$ be a UFD. The Picard group of $\mathbf{P}^n_R$
is $\mathbf{Z}$. More precisely, there is an isomorphism
$$
\mathbf{Z} \longrightarrow \Pic(\mathbf{P}^n_R),\quad
m \longmapsto \mathcal{O}_{\mathbf{P}^n_R}(m)
$$
In particular, the Picard group of $\mathbf{P}^n_k$ of projective
space over a field $k$ is $\mathbf{Z}$.
\end{lemma}

\begin{proof}
By Lemma \ref{lemma-open-subscheme-UFD} the Picard groups of the opens
$D_+(T_i) \cong \mathbf{A}^n_R$ are trivial.
Thus if $\mathcal{L}$ is an invertible module on
$\mathbf{P}^n_R$, then it is given by a $1$-cocycle with values in
the sheaf of invertible functors $\mathcal{O}^*$ for the open covering
$$
\mathbf{P}^n_R = \bigcup\nolimits_{i = 0, \ldots, n} D_+(T_i)
$$
Observe that for $i \not = j$ we have
$$
\mathcal{O}^*(D_+(T_i) \cap D_+(T_j)) =
\mathcal{O}^*(D_+(T_iT_j)) = \left(R[T_0, \ldots, T_n]_{(T_iT_j)}\right)^* =
R^* \times (T_i/T_j)^\mathbf{Z}
$$
Thus such a cocycle $(g_{ij})$ is given by units
$$
g_{ij} = u_{ij} (T_i/T_j)^{e_{ij}}
$$
with $u_{ij} \in R^*$ and $e_{ij} \in \mathbf{Z}$
satisfying the cocycle condition. The cocycle condition
over $D_+(T_iT_jT_k)$ for $\#\{i, j, k\} = 3$ tell us that
$$
u_{ik} = u_{ij} u_{jk} \quad\text{and}\quad
e_{ik} = e_{ij} = e_{jk}
$$
Whence $u_{ij} = u_{i1} u_{j1}^{-1}$ is a boundary.
Thus all isomorphism classes of invertible modules
are given by taking the cocycle with
$$
g_{ij} = (T_i/T_j)^e
$$
for some $e \in \mathbf{Z}$. Since $\mathcal{O}(n)$ has trivializing
section $T_i^n$ over $D_+(T_i)$ we see that the corresponding
cocycle of $\mathcal{O}(n)$ is $(T_i/T_j)^n$ and the proof is complete.
\end{proof}






